% GENERATED FILE. Edit canonical lessons, then run npm run generate:books.
% canonical-source-hash: fnv1a64:42f5acf8c1860845
% canonical-lessons: KA-C320-varadigara, KA-C320-sambhavisu, KA-C320-prasara, KA-C320-bittarisu, KA-C320-sirsike

\chapter{The Reporter on the Spot}
\label{ch:ka-the-reporter-on-the-spot}
\hlchaptermodality{320}

\begin{chapteropening}
\textbf{By the end of this chapter:} \emph{I can say a reporter, to occur, a broadcast, to broadcast and a headline.}

Complete the last lesson of chapter 320: I can say a reporter, to occur, a broadcast, to broadcast and a headline.
\end{chapteropening}

\section[\texorpdfstring{\kn{ವರದಿಗಾರ} (varadigāra)}{ವರದಿಗಾರ (varadigāra)}]{\kn{ವರದಿಗಾರ} (varadigāra) — a reporter}

\label{lesson:KA-C320-varadigara}



\begin{quote}
Before the new one: say the Kannada for a thief, then the Kannada for to run away.
\end{quote}

\subsection*{You'll want to know: \kn{ವರದಿಗಾರ}}
\textbf{\kn{ವರದಿಗಾರ}} — \emph{varadigāra} — \textquotedblleft{}a reporter\textquotedblright{}.

\begin{quote}
\kn{ಆ} \kn{ವರದಿಗಾರ} \kn{ಅಲ್ಲಿ} \kn{ಇದ್ದಿದ್ದರೆ} \kn{ಎಲ್ಲರಿಗೂ} \kn{ಸುದ್ದಿ} \kn{ಗೊತ್ತಾಗುತ್ತಿತ್ತು}.
\end{quote}

\begin{quote}
\emph{ā varadigāra alli iddiddare ellarigū suddi gottāguttittu.}
\end{quote}

\begin{quote}
If that reporter had been there, everyone would have known the news.
\end{quote}

Plural: \kn{ವರದಿಗಾರರು}, \emph{varadigāraru}.

\kn{ವರದಿ} is a report; -\kn{ಗಾರ} makes \textquotedblleft{}one who does it\textquotedblright{}, as in \kn{ಕೆಲಸಗಾರ}, a worker.

\begin{cousinweb}
\textbf{\kn{ವರದಿಗಾರ}} (\emph{varadigāra}) is a Kannada word built from \kn{ವರದಿ}, \textquotedblleft{}a report\textquotedblright{}, and the ending -\emph{gāra}, \textquotedblleft{}one who does\textquotedblright{}, as in \kn{ಕೆಲಸಗಾರ}. The origin of \kn{ವರದಿ} itself is not settled.
\end{cousinweb}

\begin{grammarlens}[title={What you've built}]
One more word for telling what happened.
\end{grammarlens}

\subsection*{Your turn}
\begin{itemize}
\raggedright
  \item \emph{Say it:} \emph{varadigāra}
  \item \emph{Say it:} \emph{varadigāra}, once more
  \item \emph{Say it:} \emph{ā varadigāra alli iddiddare ellarigū suddi gottāguttittu.} — If that reporter had been there, everyone would have known the news.
  \item \emph{Recall:} say the Kannada for a thief, then the Kannada for to run away, then say \emph{varadigāra} again
\end{itemize}

\begin{tcolorbox}[breakable,colback=teal!4,colframe=teal!35!black,title={Before you move on}]
What does \kn{ವರದಿಗಾರ} mean? (\textquotedblleft{}A reporter\textquotedblright{}.) Say it once more.
\end{tcolorbox}
\section[\texorpdfstring{\kn{ಸಂಭವಿಸು} (sambhavisu)}{ಸಂಭವಿಸು (sambhavisu)}]{\kn{ಸಂಭವಿಸು} (sambhavisu) — to occur; to happen}

\label{lesson:KA-C320-sambhavisu}



\begin{quote}
Before the new one: say the Kannada for to run away, then the Kannada for a reporter.
\end{quote}

\subsection*{You'll want to know: \kn{ಸಂಭವಿಸು}}
\textbf{\kn{ಸಂಭವಿಸು}} — \emph{sambhavisu} — \textquotedblleft{}to occur; to happen\textquotedblright{}.

\begin{quote}
\kn{ನಿನ್ನೆ} \kn{ರಾತ್ರಿ} \kn{ನಮ್ಮ} \kn{ಬೀದಿಯಲ್ಲಿ} \kn{ಅಪಘಾತ} \kn{ಸಂಭವಿಸಿತು}.
\end{quote}

\begin{quote}
\emph{ninne rātri namma bīdiyalli apaghāta sambhavisitu.}
\end{quote}

\begin{quote}
Last night an accident happened in our street.
\end{quote}

Past (it …): \kn{ಸಂಭವಿಸಿತು}, \emph{sambhavisitu}.

An event is the subject, so the past is \kn{ಸಂಭವಿಸಿತು}, it happened. It suits news reports.

\begin{cousinweb}
\textbf{\kn{ಸಂಭವಿಸು}} (\emph{sambhavisu}) is built on Sanskrit \emph{sambhava}, \textquotedblleft{}being or coming together, arising\textquotedblright{} (Monier-Williams), with the Kannada ending -\emph{isu} that makes verbs of Sanskrit nouns.
\end{cousinweb}

\begin{grammarlens}[title={What you've built}]
One more word for telling what happened.
\end{grammarlens}

\subsection*{Your turn}
\begin{itemize}
\raggedright
  \item \emph{Say it:} \emph{sambhavisu}
  \item \emph{Say it:} \emph{sambhavisu}, once more
  \item \emph{Say it:} \emph{ninne rātri namma bīdiyalli apaghāta sambhavisitu.} — Last night an accident happened in our street.
  \item \emph{Recall:} say the Kannada for to run away, then the Kannada for a reporter, then say \emph{sambhavisu} again
\end{itemize}

\begin{tcolorbox}[breakable,colback=teal!4,colframe=teal!35!black,title={Before you move on}]
What does \kn{ಸಂಭವಿಸು} mean? (\textquotedblleft{}To occur; to happen\textquotedblright{}.) Say it once more.
\end{tcolorbox}
\section[\texorpdfstring{\kn{ಪ್ರಸಾರ} (prasāra)}{ಪ್ರಸಾರ (prasāra)}]{\kn{ಪ್ರಸಾರ} (prasāra) — a broadcast; spreading}

\label{lesson:KA-C320-prasara}



\begin{quote}
Before the new one: say the Kannada for a reporter, then the Kannada for to occur.
\end{quote}

\subsection*{You'll want to know: \kn{ಪ್ರಸಾರ}}
\textbf{\kn{ಪ್ರಸಾರ}} — \emph{prasāra} — \textquotedblleft{}a broadcast; spreading\textquotedblright{}.

\begin{quote}
\kn{ಆ} \kn{ಪಂದ್ಯದ} \kn{ಪ್ರಸಾರ} \kn{ಸಂಜೆ} \kn{ಆರು} \kn{ಗಂಟೆಗೆ} \kn{ಶುರುವಾಯಿತು}.
\end{quote}

\begin{quote}
\emph{ā pandyada prasāra sañje āru gaṇṭege śuruvāyitu.}
\end{quote}

\begin{quote}
The broadcast of that match began at six in the evening.
\end{quote}

Radio and television both have a \kn{ಪ್ರಸಾರ}. \kn{ಪ್ರಸಾರವಾಗು} means to be broadcast.

\begin{cousinweb}
\textbf{\kn{ಪ್ರಸಾರ}} (\emph{prasāra}) is borrowed from Sanskrit \emph{prasāra}, \textquotedblleft{}spreading out, extension\textquotedblright{} (CDIAL 8835). A broadcast spreads a programme far and wide.
\end{cousinweb}

\begin{grammarlens}[title={What you've built}]
One more word for telling what happened.
\end{grammarlens}

\subsection*{Your turn}
\begin{itemize}
\raggedright
  \item \emph{Say it:} \emph{prasāra}
  \item \emph{Say it:} \emph{prasāra}, once more
  \item \emph{Say it:} \emph{ā pandyada prasāra sañje āru gaṇṭege śuruvāyitu.} — The broadcast of that match began at six in the evening.
  \item \emph{Recall:} say the Kannada for a reporter, then the Kannada for to occur, then say \emph{prasāra} again
\end{itemize}

\begin{tcolorbox}[breakable,colback=teal!4,colframe=teal!35!black,title={Before you move on}]
What does \kn{ಪ್ರಸಾರ} mean? (\textquotedblleft{}A broadcast; spreading\textquotedblright{}.) Say it once more.
\end{tcolorbox}
\section[\texorpdfstring{\kn{ಬಿತ್ತರಿಸು} (bittarisu)}{ಬಿತ್ತರಿಸು (bittarisu)}]{\kn{ಬಿತ್ತರಿಸು} (bittarisu) — to broadcast; to air}

\label{lesson:KA-C320-bittarisu}



\begin{quote}
Before the new one: say the Kannada for to occur, then the Kannada for a broadcast.
\end{quote}

\subsection*{You'll want to know: \kn{ಬಿತ್ತರಿಸು}}
\textbf{\kn{ಬಿತ್ತರಿಸು}} — \emph{bittarisu} — \textquotedblleft{}to broadcast; to air\textquotedblright{}.

\begin{quote}
\kn{ಟೀವಿ} \kn{ಆ} \kn{ಸಂದರ್ಶನವನ್ನು} \kn{ರಾತ್ರಿ} \kn{ಬಿತ್ತರಿಸಿತು}.
\end{quote}

\begin{quote}
\emph{ṭīvi ā sandarśanavannu rātri bittarisitu.}
\end{quote}

\begin{quote}
The television aired that interview at night.
\end{quote}

Past (I …): \kn{ಬಿತ್ತರಿಸಿದೆ}, \emph{bittariside}.

What is aired takes -\kn{ಅನ್ನು}. \kn{ಪ್ರಸಾರ} \kn{ಮಾಡು} says the same thing.

\begin{cousinweb}
\textbf{\kn{ಬಿತ್ತರಿಸು}} (\emph{bittarisu}) is built on \kn{ಬಿತ್ತರ}, the Kannada form of Sanskrit \emph{vistara}, \textquotedblleft{}spreading out, extent\textquotedblright{} (Monier-Williams), with the verb ending -\emph{isu}. Sanskrit \emph{v} often turns into Kannada \emph{b}.
\end{cousinweb}

\begin{grammarlens}[title={What you've built}]
One more word for telling what happened.
\end{grammarlens}

\subsection*{Your turn}
\begin{itemize}
\raggedright
  \item \emph{Say it:} \emph{bittarisu}
  \item \emph{Say it:} \emph{bittarisu}, once more
  \item \emph{Say it:} \emph{ṭīvi ā sandarśanavannu rātri bittarisitu.} — The television aired that interview at night.
  \item \emph{Recall:} say the Kannada for to occur, then the Kannada for a broadcast, then say \emph{bittarisu} again
\end{itemize}

\begin{tcolorbox}[breakable,colback=teal!4,colframe=teal!35!black,title={Before you move on}]
What does \kn{ಬಿತ್ತರಿಸು} mean? (\textquotedblleft{}To broadcast; to air\textquotedblright{}.) Say it once more.
\end{tcolorbox}
\section[\texorpdfstring{\kn{ಶೀರ್ಷಿಕೆ} (śīrṣike)}{ಶೀರ್ಷಿಕೆ (śīrṣike)}]{\kn{ಶೀರ್ಷಿಕೆ} (śīrṣike) — a headline; a title}

\label{lesson:KA-C320-sirsike}



\begin{quote}
Before the new one: say the Kannada for a broadcast, then the Kannada for to broadcast.
\end{quote}

\subsection*{You'll want to know: \kn{ಶೀರ್ಷಿಕೆ}}
\textbf{\kn{ಶೀರ್ಷಿಕೆ}} — \emph{śīrṣike} — \textquotedblleft{}a headline; a title\textquotedblright{}.

\begin{quote}
\kn{ಇವತ್ತಿನ} \kn{ಪತ್ರಿಕೆಯ} \kn{ಶೀರ್ಷಿಕೆ} \kn{ಆ} \kn{ಯುದ್ಧದ} \kn{ಬಗ್ಗೆ} \kn{ಇತ್ತು}.
\end{quote}

\begin{quote}
\emph{ivattina patrikeya śīrṣike ā yuddhada bagge ittu.}
\end{quote}

\begin{quote}
Today's newspaper headline was about that war.
\end{quote}

Plural: \kn{ಶೀರ್ಷಿಕೆಗಳು}, \emph{śīrṣikegaḷu}.

A book, a film and a song have a \kn{ಶೀರ್ಷಿಕೆ} too: its title.

\begin{cousinweb}
\textbf{\kn{ಶೀರ್ಷಿಕೆ}} (\emph{śīrṣike}) is borrowed from Sanskrit \emph{śīrṣaka}, built on \emph{śīrṣa}, \textquotedblleft{}head, top\textquotedblright{} (CDIAL 12497). A headline sits at the head of a story.
\end{cousinweb}

\begin{grammarlens}[title={What you've built}]
One more word for telling what happened.
\end{grammarlens}

\subsection*{Your turn}
\begin{itemize}
\raggedright
  \item \emph{Say it:} \emph{śīrṣike}
  \item \emph{Say it:} \emph{śīrṣike}, once more
  \item \emph{Say it:} \emph{ivattina patrikeya śīrṣike ā yuddhada bagge ittu.} — Today's newspaper headline was about that war.
  \item \emph{Recall:} say the Kannada for a broadcast, then the Kannada for to broadcast, then say \emph{śīrṣike} again
\end{itemize}

\begin{tcolorbox}[breakable,colback=teal!4,colframe=teal!35!black,title={Before you move on}]
What does \kn{ಶೀರ್ಷಿಕೆ} mean? (\textquotedblleft{}A headline; a title\textquotedblright{}.) Say it once more.
\end{tcolorbox}
